\section{Reduction to a faithful central state}\label{sec:reductions}

We first remove the obstruction to using a normal conditional expectation:
the original masa need not be the range of one.  A normal positive
functional $\rho$ on $M$ is \emph{$A$-central} if
\[
 \rho(bz)=\rho(zb)\qquad(b\in A,\ z\in M).
\]
The supremum $p$ of their support projections lies in $A$, as proved in
the corner assembly below. Given finite vector tests, that assembly chooses
one support $e\leq p$ so that $p-e$ acts arbitrarily small on those vectors.
It uses a faithful central state on $eMe$, a direct randomization on
$(1-p)M(1-p)$, and discards $p-e$. The projections $p$ and $e$ need not be
central in $M$; disjoint palettes make the pinching remove all cross-corner
terms, so the numbers of colors add. We then reduce the faithful-state case
to separable predual. All these reductions retain a bound on the number of
projections that is independent of the prescribed strong neighborhood.

\subsection{The part without central normal functionals}

For the first reduction we use Haagerup's $L^p$ spaces for an arbitrary
von Neumann algebra, constructed with a normal semifinite faithful
weight; see \cite[Section~1]{Haa79}.  If $\varphi$ is a normal positive
functional, its density $h_\varphi\in L^1(M)_+$ represents it under the
isometric identification $L^1(M)=M_*$.  The pairing is denoted by $\Tr$.
The space $L^2(M)$, with left and right multiplication and its positive
cone, is a standard form of $M$.  In particular its cone vector for
$\varphi$ is $h_\varphi^{1/2}$.  We use H\"older's inequality, cyclicity
of the $L^p$--$L^{p'}$ pairing, and polar decomposition, whose partial
isometry belongs to $M$.

The positive-cone estimate below is Haagerup's standard-form inequality
\cite[Lemma~2.10]{Haa75}; see also the earlier work of Powers and
St{\o}rmer \cite[Lemma~4.1]{PS70} and Araki
\cite[Theorem~4(8), equation~(5.10), and Section~8, Remark~2]{Ara74}. We include its proof
so that no faithfulness assumption on the two functionals is implicit.
For normal positive $\varphi,\psi$,
\begin{equation}\label{eq:positive-cone-estimate}
 \|h_\varphi^{1/2}-h_\psi^{1/2}\|_2^2
 \leq \|\varphi-\psi\|.
\end{equation}
Indeed, put $\xi_1=h_\varphi^{1/2}$, $\xi_2=h_\psi^{1/2}$ and
$D=\xi_1-\xi_2$.  Its sign $s\in M$ satisfies
$\|s\|\leq1$ and $sD=Ds=|D|$.  Writing $D=D_+-D_-$, cyclicity gives
\begin{align*}
 \Tr\bigl(s(h_\varphi-h_\psi)\bigr)
 &=\Tr\bigl(s(D\xi_1+\xi_2D)\bigr)
   =\Tr\bigl(|D|(\xi_1+\xi_2)\bigr)\\
 &=\|D\|_2^2+2\Tr(D_-\xi_1)+2\Tr(D_+\xi_2)
 \geq\|D\|_2^2.
\end{align*}
The last two pairings are nonnegative by self-duality of the positive
cone, and the first expression is at most $\|\varphi-\psi\|$.

\begin{lemma}\label{lem:no-central}
Let $A\subseteq M$ be a masa such that $M$ has no nonzero $A$-central
normal positive functional.  Let $x=x^*\in M$ satisfy $\|x\|\leq1$,
and let $0<\eps<1$.  Put $r_{\rm p}=\lceil4/\eps\rceil$.
For every normal state $\varphi$ on $M$ and every $\rho>0$ there are a
partition $p_1,\ldots,p_{r_{\rm p}}$ in $A$, an $a=a^*\in A$ with
$\|a\|\leq1$, and a projection $q\in M$ such that
\[
 \varphi(1-q)<\rho,
 \qquad
 \left\|q\left(\sum_{i=1}^{r_{\rm p}}p_i x p_i-a\right)q\right\|
 \leq\eps.
\]
\end{lemma}

\begin{proof}
Write $\xi=h_\varphi^{1/2}$ and introduce the symmetric map
\[
 T(z)=h_\varphi^{1/4}z h_\varphi^{1/4}\colon M\longrightarrow L^2(M).
\]
It is bounded and takes weak-star convergence to weak convergence:
testing against $\zeta\in L^2(M)$ uses the $L^1$ element
$h_\varphi^{1/4}\zeta^*h_\varphi^{1/4}$.
For a finite partition $P$ in $A$, the pinching
$S_P=\sum_{e\in P}e(\,\cdot\,)e$ also acts on $L^2(M)$.
There it is an orthogonal projection.  These projections decrease under
refinement to the projection onto the $A$-central vectors.  Such a vector
must be zero: otherwise its left vector functional would be a nonzero
$A$-central normal positive functional.  To see centrality, move a unitary
of $A$ from the left of the vector to its right, where it commutes with
the left action of $M$.  Consequently
\[
 \|S_P(\xi)\|_2\longrightarrow0.
\]
For $z\in M$, cyclicity and H\"older's inequality give
\begin{equation}\label{eq:symmetric-block-sum}
 \sum_{e,f\in P}\|T(ezf)\|_2^2
 =\Tr\bigl(z^*S_P(\xi)zS_P(\xi)\bigr)
 \leq\|z\|^2\|S_P(\xi)\|_2^2.
\end{equation}
In particular $\|T(eze)\|_2^2\leq\|e\xi e\|_2^2$ when $\|z\|\leq1$.

Fix $\kappa>0$ and a coarse partition $P_0$ with
$\|S_{P_0}(\xi)\|_2<\kappa$.
A weak-star convergent subnet of $S_P(x)$, as $P$ refines $P_0$, has a
self-adjoint contraction limit $b\in A$: the limit commutes with every
partition in $A$, hence belongs to the masa.  Since the weak and norm
closures of a convex subset of a Hilbert space agree, there are
refinements $P^1,\ldots,P^d$ of $P_0$ and probabilities $t_1,\ldots,t_d$
such that
\[
 \left\|\sum_{j=1}^d t_jT(S_{P^j}(x))-T(b)\right\|_2<\kappa.
\]
For each $e\in P_0$, independently choose $j_e$ with this distribution
and use the partition $P^{j_e}$ inside $e$.  The resulting partition $R$
refines $P_0$.  The mean of $T(S_R(x))$ is the displayed convex
combination.  Independence and Equation~\eqref{eq:symmetric-block-sum}
bound its variance by
\[
 \sum_{e\in P_0}\|e\xi e\|_2^2
 =\|S_{P_0}(\xi)\|_2^2<\kappa^2.
\]
Thus one realization satisfies
$\|T(S_R(x)-b)\|_2<\sqrt2\kappa$.

Now color the pieces of this fixed $R$ independently and uniformly with
$r_{\rm p}$ colors, and let $S_{\rm col}$ be the resulting color
pinching.  Its mean is
\[
 \mathbb E S_{\rm col}(x)
 =\frac{x}{r_{\rm p}}+\left(1-\frac1{r_{\rm p}}\right)S_R(x).
\]
The indicators of equal colors for distinct unordered pairs of pieces
are uncorrelated, even when the pairs share a piece.  Applying
$\|u+v\|_2^2\leq2\|u\|_2^2+2\|v\|_2^2$ and
Equation~\eqref{eq:symmetric-block-sum}, the variance after applying $T$
is at most $2\|S_R(\xi)\|_2^2<2\kappa^2$.
Therefore a coloring can be chosen so that the self-adjoint element
\[
 y=S_{\rm col}(x)-\left(1-\frac1{r_{\rm p}}\right)b
                     -\frac{x}{r_{\rm p}}
\]
satisfies $\|T(y)\|_2<2\sqrt2\kappa$.
Also $\|y\|\leq2$.  Since $\kappa$ is arbitrary, this gives arbitrarily
small symmetric seminorm with the number of colors already fixed.

It remains to convert this seminorm into a large compression.  Set
$t=\|T(y)\|_2$ and $\psi(z)=\varphi(y^*zy)$.
Small $t$ need not control $\|y\|_\varphi$, so we first seek $q_0$ with
both $\varphi(1-q_0)$ and $\psi(q_0)$ small; these bounds will control
$\|q_0yq_0\|_\varphi$ for the spectral cut.
Polar decomposition gives
$h_\psi=yh_\varphi y^*$ and
$h_\psi^{1/2}=y\xi u^*$ for a partial isometry $u\in M$.
Hence
\[
 0\leq\Tr(\xi h_\psi^{1/2})
 \leq\|\xi y\xi\|_1
 =\|h_\varphi^{1/4}T(y)h_\varphi^{1/4}\|_1
 \leq t.
\]
Here the first inequality uses the positive cone, and the last uses
$\|h_\varphi^{1/4}\|_4=1$.
Let $q_0$ support the positive part in the Jordan decomposition of
$\varphi-\psi$.  Equation~\eqref{eq:positive-cone-estimate} implies
\begin{equation}\label{eq:jordan-overlap}
 \varphi(1-q_0)+\psi(q_0)
 =\frac{\varphi(1)+\psi(1)-\|\varphi-\psi\|}{2}
 \leq t.
\end{equation}
Moreover
\[
 \|q_0yq_0\xi\|_2
 \leq\|q_0y\xi\|_2+\|q_0y(1-q_0)\xi\|_2
 \leq\psi(q_0)^{1/2}+2\varphi(1-q_0)^{1/2}
 \leq3\sqrt t.
\]
Inside $q_0Mq_0$, let $q$ be the spectral projection of $q_0yq_0$ for
$[-\eps/2,\eps/2]$.  The spectral Markov inequality gives
\begin{equation}\label{eq:symmetric-large-projection}
 \varphi(1-q)
 \leq\varphi(1-q_0)+\frac4{\eps^2}\varphi((q_0yq_0)^2)
 \leq\left(1+\frac{36}{\eps^2}\right)t.
\end{equation}
Taking $\kappa$ sufficiently small makes this less than $\rho$.
With $a=(1-1/r_{\rm p})b$, we have $\|a\|\leq1$ and
\[
 \|q(S_{\rm col}(x)-a)q\|
 \leq\frac\eps2+\frac1{r_{\rm p}}
 \leq\frac{3\eps}{4}.
\]
This proves the lemma.
\end{proof}

\subsection{Assembly of corners}

\begin{proposition}\label{prop:state-reduction}
Fix $0<\eps<1$ and an integer $R\geq1$.  Suppose that for every masa
$A\subseteq M$ admitting a faithful $A$-central normal state $\varphi$,
every self-adjoint contraction $x\in M$, and every $\rho>0$, there exist
a partition of $1$ into $R$ projections of $A$, a self-adjoint
$a\in A$ of norm at most one, and a projection $q\in M$ with
$\varphi(1-q)<\rho$ and
\[
 \left\|q\left(\sum_{i=1}^R p_i x p_i-a\right)q\right\|\leq\eps.
\]
Then every masa in every represented von Neumann algebra has the
strong-operator paving property of Theorem~\ref{thm:main} at error
$\eps$, with
\[
 r=R+\lceil4/\eps\rceil+1.
\]
In particular this number is chosen before the finite vector set and
the strong-neighborhood tolerance.
\end{proposition}

\begin{proof}
Let $p$ be the supremum of the supports of the $A$-central normal
positive functionals on $M$.  Each support belongs to $A$, since
conjugation by every unitary of $A$ preserves the functional and its
support.  These supports form a directed family: the support of a sum
is the join of the supports.  The corner $(1-p)M(1-p)$ has no nonzero
$A(1-p)$-central normal positive functional, since compression would
extend one to an $A$-central functional on $M$.

Given a finite $F\subset H$ and $\delta>0$, set
\[
 \chi(z)=\sum_{\xi\in F}\langle z\xi,\xi\rangle.
\]
Normality supplies a support $e\leq p$ in the directed family with
$\chi(p-e)$ as small as desired.  On the nonzero corner $eMe$, the
corresponding functional restricts, after normalization, to a faithful
normal state central on $Ae$.

For any faithful normal state $\omega$, small $\omega$-mass on
projections implies uniformly small mass for each fixed normal positive
functional.  Indeed, a contrary sequence in the positive unit ball
would have a weak-star convergent subnet whose limit has zero
$\omega$-mass but positive mass for that functional.  Faithfulness
forces the positive limit to be zero, a contradiction.
Apply the assumed result on $eMe$, using this observation to make the
lost $\chi$-mass small.  On $(1-p)M(1-p)$ apply
Lemma~\ref{lem:no-central} to the normalized restriction of $\chi$ if
it is nonzero; otherwise take $q=0$ on this corner.  Empty corners
require no construction.

Use disjoint palettes for the two corner partitions, and the additional
projection $p-e$ as the last piece.  Take corner sums of the two
diagonal terms and of the two large projections, with zero on $p-e$.
Every pinching piece is supported in one of the three corners, so all
off-corner terms of $x$ disappear.  Thus the compressed operator norm
is at most $\eps$, and the diagonal norm is at most one.  The choices
can ensure $\chi(1-q)<\delta^2$, which gives
\[
 \|(1-q)\xi\|^2\leq\chi(1-q)<\delta^2\qquad(\xi\in F).
\]
Zero projections pad either palette when necessary.  The total number
is the stated $r$, independently of $F$ and $\delta$.  Scaling handles
an arbitrary nonzero self-adjoint $x$; for $x=0$ the assertion is
immediate.
\end{proof}

\subsection{State norm and separable reduction}

We now work with a faithful $A$-central normal state $\varphi$.
Write
\[
 \|z\|_\varphi=\varphi(z^*z)^{1/2},
 \qquad
 M_\varphi=\{b\in M:\varphi(bz)=\varphi(zb)\text{ for all }z\in M\}.
\]
The centralizer $M_\varphi$ is the fixed algebra of the modular group
$\sigma^\varphi$.  In particular $A\subseteq M_\varphi$.
We first obtain approximation by diagonal pinchings in the state norm,
then use it to construct a separable inclusion containing the operator.

\begin{lemma}\label{lem:pinching}
Let $A\subseteq M$ be a masa and $\varphi$ a faithful $A$-central normal
state.  There is a faithful normal $\varphi$-preserving conditional
expectation $E_A\colon M\to A$.  On bounded subsets of $M$, convergence
to zero in $\|\cdot\|_\varphi$ is equivalent to strong convergence in the
faithful GNS representation.  For every $z\in M$,
\begin{equation}\label{eq:pinching-limit}
 \|S_P(z)-E_A(z)\|_\varphi\longrightarrow0
\end{equation}
as the finite partitions $P$ of $A$ refine.  The same conclusion holds
along any increasing sequence of finite partitions generating $A$.
\end{lemma}

\begin{proof}
Takesaki's Conditional Expectation Theorem \cite{Tak72} applies because
the modular group fixes $A$ pointwise; both the state and its restriction
are faithful, normal, and finite.  Let $\Omega$ be the GNS state vector.
It is separating and therefore cyclic for the commutant.  If a bounded
net $z_j$ satisfies $z_j\Omega\to0$, then
$z_jc'\Omega=c'z_j\Omega\to0$ for every commutant element $c'$.
Density and boundedness give strong convergence.  The converse follows
by testing on $\Omega$.

Since the partition projections lie in the centralizer, $S_P$ induces
an orthogonal projection on the GNS space; equivalently, it is the
state-preserving conditional expectation onto the commutant in $M$ of
the projections of $P$.  These Hilbert-space projections decrease under
refinement.  On the vector $z\Omega$, let their limit be $\eta$.
A bounded weak-star cluster point $b$ of $S_P(z)$ commutes with every
partition of $A$ and hence lies in $A$.  Its state vector is $\eta$.
Moreover $E_A S_P(z)=E_A(z)$, so $b=E_A(z)$.  This proves
Equation~\eqref{eq:pinching-limit}.  If an increasing sequence generates
$A$, its limit commutes with the generating projections, and the same
argument applies.
\end{proof}

\begin{lemma}\label{lem:separable-reduction}
Let $A\subseteq M$ be a masa with a faithful $A$-central normal state
$\varphi$, and let $x=x^*\in M$.  There exist unital von Neumann
subalgebras $A_0\subseteq A$ and $M_0\subseteq M$ such that $x\in M_0$,
$M_0$ has separable predual, $A_0$ is a masa in $M_0$, and
$E_A(M_0)\subseteq A_0$.  The restriction of $\varphi$ is faithful,
normal, and $A_0$-central.  Thus a uniform faithful-state paving result
for inclusions with separable predual implies the same result without
that restriction.
\end{lemma}

\begin{proof}
Start with $D_1=W^*(1,x)$.  Inductively, choose a countable
$\|\cdot\|_\varphi$-dense subset $\mathcal D_j$ of the unit ball of
$D_j$.  Such a set exists for every countably generated $D_j$: the
separable C*-algebra of its generators has dense state vectors by
Kaplansky's Density Theorem.  Its GNS representation is faithful and
normal on a separable Hilbert space, so $D_j$ has separable predual.
Let $D_{j+1}$ be generated by $D_j$, all $E_A(z)$ for
$z\in\mathcal D_j$, and the projections of finite partitions
$P(z,k)\subset A$ satisfying
\[
 \|S_{P(z,k)}(z)-E_A(z)\|_\varphi<1/k
 \qquad(z\in\mathcal D_j,\ k\geq1).
\]
Lemma~\ref{lem:pinching} supplies these partitions.  Only countably
many generators have been added.

Put $M_0=W^*(\bigcup_jD_j)$ and let $A_0$ be generated unitally by all
the adjoined elements of $A$.  Bounded Kaplansky approximation by the
union algebra, followed by approximation from the sets $\mathcal D_j$,
shows that their union is state-norm dense in the unit ball of $M_0$.
Here one may first approximate in the norm closure of the union and
then make a norm approximation and a harmless rescaling to preserve
the bound.  Since $E_A$ is state-norm contractive, every $E_A(z)$ for
$z\in M_0$ is a bounded state-norm limit of elements of $A_0$.
Lemma~\ref{lem:pinching} and strong closedness imply $E_A(z)\in A_0$.

If a unit-ball element $z\in M_0$ commutes with $A_0$, each selected
partition satisfies $S_P(z)=z$.  For $z'\in\mathcal D_j$ and its
selected partition $P=P(z',k)$, contractivity gives
\[
 \|z-E_A(z)\|_\varphi
 \leq 2\|z-z'\|_\varphi
       +\|S_P(z')-E_A(z')\|_\varphi.
\]
The right side can be made arbitrarily small; faithfulness gives
$z=E_A(z)\in A_0$.  Thus $A_0$ is maximal abelian in $M_0$.
The algebra $M_0$ is countably generated and has separable predual by
the same GNS argument.  All the assertions about the restricted state
are immediate. Any paving produced in $M_0$ has the same operator-norm
bound and the same lost $\varphi$-mass when regarded in $M$. This is the
faithful-state estimate required by Proposition~\ref{prop:state-reduction}.
\end{proof}
